\input{preamble-template.tex}
\hsslug{pd-trigger-lite}
\begin{document}

\hstitleblock{Board guide \textperiodcentered\ September 2026}{pd-trigger-lite}%
{A 25 by 15 mm USB-C PD trigger: plug in a PD charger and get 9, 12, 15 or 20 V at up to 3 A on two solder pads.}

This board asks a USB-C Power Delivery charger for one fixed voltage and passes the charger's
VBUS straight to two 2.54 mm output holes. There is no regulator and no switch in the power path,
so the output is exactly what the charger delivers. The voltage is picked by one small solder
link, and the board ships set to 12 V. It is meant for bench use: powering a strip, a router or a
prototype from a laptop charger.

\hsplate{top.png}{Top view}{25 x 15 mm, 2 layers, 1 oz}{render}%
{USB-C receptacle on the left, the voltage links along the top, the PD controller at the bottom,
output pads on the right (+ is VBUS, - is ground). The red LED sits beside the output pads.}

\newpage
\section{Schematic}
\hslead{One controller, one capacitor, four resistors and a link. Everything else is copper.}

The receptacle J1 is a 6-pin power-only part: two VBUS contacts, two ground contacts and the two
CC lines. The controller U1 (WCH CH224A) talks to the charger over CC1 and CC2 and needs nothing
else to negotiate. It is powered from VBUS through the 0 ohm link R7, which only exists so the
controller's pins can be reached with thin copper, and C1 (1 uF, 50 V) decouples it. VBUS runs as
a wide copper pour from J1 to the output pad J2.

U1 reads the resistance from its CFG1 pin to ground once, at power-up, and requests the matching
voltage: 6.8k gives 9 V, 24k gives 12 V, 56k gives 15 V and 120k gives 20 V. Each of those
resistors (R1 to R4) reaches ground only through its own link, so the one closed link decides the
voltage. R6 and D1 light the LED from VBUS.

\hsplate{schematic.png}{Schematic}{kicad/pd-trigger-lite.kicad\_sch}{ERC 0 errors, 0 warnings}%
{Nets join by label name. JP1 to JP3 are empty 0603 lands; R5 is the 0 ohm resistor fitted in the
12 V position.}

\newpage
\section{Using it}
\hslead{Set the voltage before the first plug-in, check it with a meter, then connect the load.}

\begin{hsblock}
\begin{tabularx}{\linewidth}{L{2.6cm}L{2.2cm}Y}
\hstoprule
\hshead{Voltage} & \hshead{Link} & \hshead{What is fitted} \\
\hstoprule
9 V  & JP1 & R1 6.8k in series \\ \hline
12 V & R5  & R2 24k in series; shipped fitted with a 0 ohm resistor \\ \hline
15 V & JP2 & R3 56k in series \\ \hline
20 V & JP3 & R4 120k in series \\
\hstoprule
\end{tabularx}
\end{hsblock}

To change the voltage, unplug the board, remove the 0 ohm resistor R5 with the iron, and either
solder it onto the land marked with the voltage you want or bridge that land with a blob of solder.
The four links sit in a row under the resistor row, labelled 9, 12, 15 and 20 on the silkscreen.

\begin{hscallout}{Exactly one link}
With no link closed, CFG1 floats and the request is undefined. With two closed, two resistors sit
in parallel and the controller reads a value that is in no table. Close one link, never zero and
never two.
\end{hscallout}

\textbf{Output.} The two plated holes take 20 AWG wire or a 1x2 2.54 mm header, which you solder
yourself; JLC does not fit them. The rating is 3 A at any voltage. A charger without an e-marked
cable never offers more than 3 A, and the board makes no 5 A claim.

\textbf{LED.} D1 is lit whenever VBUS is present. Its current rises with the voltage, so it glows
faintly at 5 V and clearly at 12 V and above. A dim LED after plugging in means the charger stayed
at 5 V: it did not offer the voltage you picked, or no link is closed.

\textbf{First power-up.} Plug in with nothing on the output, measure the voltage across the pads,
and only then connect the load. The charger must offer the chosen voltage as a fixed PDO; many
phone chargers stop at 9 V. The board has no fuse and no TVS diode, because the controller's
datasheet asks for neither. A hot-plug spike at 20 V is the one plausible way to damage it, so at
20 V plug the charger into the wall first and the cable into the board last.

\newpage
\section{What is on it}
\hslead{Eleven parts placed by JLC, two of them Extended. Everything else is bare copper.}

\begin{hsblock}
\begin{tabularx}{\linewidth}{L{1.9cm}Y L{2.0cm}L{1.6cm}R{0.8cm}}
\hstoprule
\hshead{Ref} & \hshead{Part} & \hshead{LCSC} & \hshead{Type} & \hshead{Qty} \\
\hstoprule
U1 & CH224A PD sink controller, ESSOP-10 & C42459160 & Extended & 1 \\ \hline
J1 & USB-C receptacle 6P, 3 A 20 V (TYPE-C-6M-001) & C2798175 & Extended & 1 \\ \hline
C1 & 1 uF 50 V X5R 0603 & C15849 & Basic & 1 \\ \hline
R1 & 6.8k 1\% 0603 & C23212 & Basic & 1 \\ \hline
R2 & 24k 1\% 0603 & C23352 & Basic & 1 \\ \hline
R3 & 56k 1\% 0603 & C23206 & Basic & 1 \\ \hline
R4 & 120k 1\% 0603 & C25808 & Basic & 1 \\ \hline
R5, R7 & 0 ohm 0603 & C21189 & Basic & 2 \\ \hline
R6 & 10k 1\% 0603 & C25804 & Basic & 1 \\ \hline
D1 & red LED 0603 & C2286 & Basic & 1 \\
\hstoprule
\end{tabularx}
\end{hsblock}

Not assembled: JP1 to JP3 are empty 0603 lands (the 9, 15 and 20 V links), and J2 is the pair of
output holes. The two Extended parts each add JLC's per-part feeder fee; no Basic part can replace
either of them.

\section{What it costs}
\hslead{About 27 USD for five assembled boards before shipping, as an estimate.}

\begin{hsstatrow}[3]
\hsstat{24.33}{USD, JLC estimate for 5 boards: PCB 2.00, setup 8.00, 2 Extended fees 6.00, stencil 8.00, joints 0.33.}
\hsstat{2.78}{USD, parts for 5 boards at LCSC unit prices, about 0.56 per board.}
\hsstat{27.1}{USD total for 5 boards, before shipping and tax.}
\end{hsstatrow}

These are estimates from the skill's price table (verified 2026-07-24), not a quote. Live JLC
prices have run 1.9 to 3.1 times higher than that table for bare PCBs, so expect the PCB line to be
nearer 4 to 6 USD. Shipping depends on the country and the carrier and is not included. The cart
at \hsurl{https://cart.jlcpcb.com/quote} is the only real price.

\hsprovenance{Estimate from fab/quote.json (order\_quote.py, qty 5, HASL, economic assembly) and
kicad/parts.json prices, read 2026-09-24.}

\section{Ordering at JLCPCB}
\hslead{Four files go up: the gerber zip, then BOM.csv and CPL.csv for assembly.}

\begin{enumerate}
\item Open the JLC quote page (link above) and upload the gerber zip,
\texttt{fab/pd-trigger-lite\_gerbers.zip}.
\item Check that the viewer reads 2 layers, 25 x 15 mm. Keep 1.6 mm, 1 oz, HASL, green, quantity 5.
\item Turn on PCB Assembly: Economic, top side, 5 boards.
\item Upload \texttt{fab/BOM.csv} and \texttt{fab/CPL.csv}. Confirm all eleven parts matched their
LCSC numbers and that JP1, JP2, JP3 and J2 are absent, since those are not assembled.
\item In the placement preview, check three parts. D1: its cathode mark must face the right
board edge. U1: the pin-1 dot sits at the corner nearest the capacitor C1. J1: the receptacle mouth
points off the left edge. No rotation correction was applied to the CPL, because every footprint
comes from JLC's own library.
\item Run JLCDFM on the gerbers as a second opinion, then pay. Payment is yours to do.
\end{enumerate}

\hsprovenance{Gates on 2026-09-24: ERC 0/0, DRC 0 errors 0 unconnected (reports/drc\_final.json),
verify and DFM pass.}
\end{document}
